% from ega4-2-fr.tex
\oldpage[IV-2]{213}

Montrons en second lieu que si $A$ vérifie $c)$, il en est de même de toute
$A$-algèbre de type fini $B$. En premier lieu, pour tout idéal premier
$\mathfrak q$ de $B$, $B_{\mathfrak q}$ est anneau local en un idéal premier
de $S^{-1}B$, où $S=A-\mathfrak p$, $\mathfrak p$ étant l'image réciproque de
$\mathfrak q$ dans $A$; comme par hypothèse tout anneau quotient intègre de
$A_{\mathfrak p}$ est un anneau japonais, il en est de même de tout anneau
quotient intègre de $B_{\mathfrak q}$, puisque $S^{-1}B$ est une
$A_{\mathfrak p}$-algèbre de type fini (7.6.5). D'autre part, si $C$ est un
quotient intègre de $B$, de corps des fractions $K$, $K'$ une extension
radicielle finie de $K$ et $C'$ une sous-$C$-algèbre finie de $K'$, de corps
des fractions $K'$, $C'$ est une $A$-algèbre intègre de type fini, et en vertu
de (6.13.7), l'hypothèse $c)$ sur $A$ entraîne que le spectre de $C'$ contient
une partie ouverte non vide dont tous les points sont normaux; autrement dit,
il y a un $g'\neq0$ dans $C'$ tel que $C'_{g'}$ soit intégralement clos, ce qui
prouve notre assertion.

Ce résultat montre que, pour prouver l'équivalence de $a)$, $b)$ et $c)$, il
suffit de montrer que $c)$ entraîne $b)$, et même de prouver que pour tout
anneau quotient intègre $B$ de $A$, de corps des fractions $K$, la clôture
intégrale $B'$ est un $B$-module de type fini. Or, la condition $c)$ entraîne
qu'il existe $f\neq0$ dans $B$ tel que $B_f$ soit intégralement clos, et on a
ainsi vérifié pour $B$ la condition (i) de (6.13.6). D'autre part, la condition
(ii) de (6.13.6) est aussi vérifiée par $B$; en effet, pour tout idéal maximal
$\mathfrak q$ de $B$, $B_{\mathfrak q}$ est quotient d'un anneau local
$A_{\mathfrak p}$, où $\mathfrak p$ est maximal dans $A$; comme
$A_{\mathfrak p}$ est par hypothèse un anneau japonais, il en est de même de
$B_{\mathfrak q}$; en outre, pour tout idéal premier $\mathfrak r$ de $B$,
$B_{\mathfrak r}$ est anneau de fractions d'un $B_{\mathfrak q}$, où
$\mathfrak q$ est un idéal maximal de $B$, donc est encore un anneau japonais,
ce qui prouve notre assertion.

On a en outre vu plus haut que si $A$ vérifie les conditions équivalentes
$a)$, $b)$, $c)$, il en est de même de toute $A$-algèbre de type fini. D'autre
part, si $A$ vérifie $b)$, il en est de même de tout anneau de fractions
$S^{-1}A$, car tout quotient intègre de $S^{-1}A$ est de la forme
$S^{-1}A/S^{-1}\mathfrak p$, où $\mathfrak p$ est un idéal premier de $A$, et
comme $A/\mathfrak p$ est par hypothèse un anneau japonais, il en est de même
de $S^{-1}(A/\mathfrak p)$ (0, 23.1.1). Cela achève de prouver la dernière
assertion de l'énoncé.
\end{proof}

\begin{corollary}[7.7.3]
\label{IV.2.s7.7.3-fr}
Soient $A$ un anneau noethérien vérifiant les conditions suivantes :
\begin{enumerate}
\item[\emph{(i)}] Pour tout idéal maximal $\mathfrak m$ de $A$, les fibres
formelles de $A_{\mathfrak m}$ sont géométriquement réduites.
\item[\emph{(ii)}] Les conditions équivalentes de (6.12.4) sont vérifiées.
\end{enumerate}
Alors $A$ est universellement japonais.
\end{corollary}

\begin{proof}
En effet, il résulte alors de (7.6.4) que tout anneau quotient intègre de
$A_{\mathfrak m}$ est un anneau japonais; d'où la conclusion, en vertu de
(7.7.2).
\end{proof}

\begin{corollary}[7.7.4]
\label{IV.2.s7.7.4-fr}
Un anneau de Dedekind dont le corps des fractions est de caractéristique 0 (en
particulier $\mathbf Z$) est un anneau universellement japonais. Un anneau
local noethérien dont les fibres formelles sont géométriquement réduites (en
particulier un anneau local noethérien complet) est universellement japonais.
\end{corollary}

\begin{proof}
La première assertion résulte de (7.7.3), (7.6.7, (i)) et (6.12.6). La seconde
résulte de (7.6.4) et (7.7.2).
\end{proof}

\subsection{Anneaux excellents}
\label{IV.2.ss7.8-fr}

\begin{env}[7.8.1]
\label{IV.2.s7.8.1-fr}
Les numéros précédents de ce paragraphe (ainsi que (6.11), (6.12) et (6.13))
ont été consacrés à l'étude de problèmes que l'on rencontre fréquemment dans
l'utilisation des anneaux et des préschémas noethériens, et qui peuvent se
grouper dans les types suivants :

